\lysh{22059_SOLUTION}{1}
\textbf{ΛΥΣΗ}

α) Το $\text{A}(1, -3)$ δεν είναι σημείο της ευθείας $\varepsilon: 4x + 6y = 1$, καθώς οι συντεταγμένες του σημείου δεν επαληθεύουν την εξίσωση της ευθείας μια και:
\[ 4 \cdot 1 + 6 \cdot (-3) = -14 \neq 1. \]

β) $4x + 6y = 1 \Rightarrow y = -\frac{2}{3}x + \frac{1}{6}$ άρα, λόγω παραλληλίας, η κλίση της ευθείας $\varepsilon$ είναι $\lambda = -\frac{2}{3}$. Οπότε, η ζητούμενη ευθεία διέρχεται από το σημείο $\text{A}(1, -3)$ και έχει κλίση $\lambda = -\frac{2}{3}$, συνεπώς έχει εξίσωση:
\[ y - (-3) = -\frac{2}{3}(x - 1) \]
ή ισοδύναμα
\[ 2x + 3y + 7 = 0. \]
\end{lysh}